% Terjemahan Bahasa Indonesia dari chapter1.tex baris 9-31.
% Sumber: Methods of Algebra, Volume 2, commit
% 9a5803ff77dd3257484cb177f851a73770a59dd3 (CC BY 4.0).
% stable-unit-id: o014.aljabr2.chapter1.overview

% segment-id: o014.li2.chapter1.overview.h001
\chapter{Pelengkap Teori Kategori}\label{sec:categories}
\hypertarget{o014.aljabr2.chapter1.overview}{ }

% segment-id: o014.li2.chapter1.overview.p001
Sesuai dengan judulnya, bab ini memperkenalkan sejumlah perangkat teori
kategori yang diperlukan dalam bab-bab berikut, agar pembaca dapat menyegarkan
kembali pengetahuannya. Sebagian materi telah dibahas dalam \cite{Li1} atau
buku ajar sejenis. Namun, untuk meninjau kembali konsep dan menyeragamkan
notasi, materi tersebut tetap disajikan ulang di sini. Untuk istilah dasar
lainnya, pembaca dapat merujuk pada Pendahuluan buku ini.

% segment-id: o014.li2.chapter1.overview.l001
\begin{itemize}
	% segment-id: o014.li2.chapter1.overview.i001
	\item Dalam \S\S~\sourcecrossref{sec:subquot}{1.1}--%
	\sourcecrossref{sec:images}{1.2}, mula-mula diperkenalkan secara ringkas
	subobjek dan objek kuosien, termasuk sifat-sifat umum monomorfisme dan
	epimorfisme. Selanjutnya diperkenalkan citra $\Image(f)$ dan kocitra
	$\Coim(f)$ suatu morfisme $f$, serta morfisme yang memenuhi
	\emph{citra sama dengan kocitra}, yang disebut morfisme ketat. Semua ini
	merupakan konsep dasar yang berlaku dalam kategori umum.

	% segment-id: o014.li2.chapter1.overview.i002
	\item \S\S~\sourcecrossref{sec:Ker-Coker}{1.3}--%
	\sourcecrossref{sec:linear-cat}{1.4} merupakan bagian linear bab ini.
	Di sana diperkenalkan jenis khusus kategori yang himpunan morfismenya
	mempunyai struktur grup abelian, yang disebut kategori-$\cate{Ab}$;
	funktor yang mempertahankan penjumlahan morfisme disebut funktor aditif.
	Dalam kategori-kategori tersebut, kernel dan kokernel suatu morfisme dapat
	didefinisikan sebagai perumuman kernel dan kokernel homomorfisme modul.
	Kategori-$\cate{Ab}$ yang mempunyai objek nol dan semua produk hingga
	disebut kategori aditif. Dalam kategori aditif dapat dibahas jumlah
	langsung hingga dari objek, sedangkan morfisme di antara jumlah-jumlah
	langsung hingga dapat diolah seperti operasi matriks.

	Lebih jauh, struktur penjumlahan dapat ditingkatkan menjadi perkalian
	skalar oleh suatu gelanggang komutatif $\Bbbk$. Dari sini diperoleh
	pengertian kategori-$\Bbbk\dcate{Mod}$ dan kategori $\Bbbk$-linear.
	Kategori-kategori berstruktur linear inilah yang menjadi perhatian utama
	buku ini.

	% segment-id: o014.li2.chapter1.overview.i003
	\item \S\S~\sourcecrossref{sec:limit-functor}{1.5}--%
	\sourcecrossref{sec:filtered-indlim}{1.6} memuat dasar-dasar tentang
	limit. Secara khusus,
	Definisi~\sourcecrossref{def:create-limit}{1.5.2} memperkenalkan istilah
	untuk menyatakan apakah suatu funktor $F: \mathcal{C} \to \mathcal{D}$
	\emph{mempertahankan} limit (dari $\mathcal{C}$ ke $\mathcal{D}$) atau
	\emph{menciptakan} limit (dari $\mathcal{D}$ ke $\mathcal{C}$).
	Selain itu dibahas pula funktor kofinal dan limit induktif terfilter;
	istilah-istilah ini praktis dan penting.

	% segment-id: o014.li2.chapter1.overview.i004
	\item Ekstensi Kan kiri dan kanan yang dibahas dalam
	\S\S~\sourcecrossref{sec:Kan-extension}{1.7}--%
	\sourcecrossref{sec:Kan-as-limit}{1.8} merupakan operasi dasar dalam
	teori kategori. Dengan syarat yang sesuai, ekstensi tersebut dapat
dibangun menggunakan $\varinjlim$ atau $\varprojlim$. Argumennya cukup
berliku, tetapi gagasan utamanya sederhana: merumuskan persoalan ekstensi
funktor sebagai suatu \emph{aproksimasi terbaik} dalam bentuk limit. Dicatat
	pula Teorema~\sourcecrossref{prop:Quillen-adjunction-gen}{1.7.7}, hasil
	yang tidak serta-merta jelas dan menerangkan hubungan antara ekstensi Kan
	kiri/kanan absolut dan pasangan adjoin.

	% segment-id: o014.li2.chapter1.overview.i005
	\item Lokalisasi Gabriel--Zisman yang diperkenalkan dalam
	\S~\sourcecrossref{sec:cat-localization}{1.9} (selanjutnya cukup disebut
	\emph{lokalisasi} dalam buku ini) merupakan konstruksi yang secara formal
	menambahkan invers bagi suatu keluarga morfisme $S$ dalam kategori.
	Konstruksi ini serupa dengan lokalisasi gelanggang, tetapi rincian
	konkretnya jauh lebih rumit. Agar hasil setelah penginversan lebih mudah
	dikendalikan, dalam praktik $S$ sering diandaikan sebagai apa yang disebut
	\emph{sistem multiplikatif} dalam
	Definisi~\sourcecrossref{def:multiplicative-family}{1.9.5}. Penerapan
	utama lokalisasi dalam buku ini ialah membangun kategori turunan; penerapan
	berikutnya ialah mendefinisikan kuosien Serre suatu kategori abelian dalam
	\S~\sourcecrossref{sec:Serre-subcat}{2.9}.
	\S~\sourcecrossref{sec:functor-localization}{1.10} akan menerangkan cara
	melakukan ekstensi Kan sepanjang lokalisasi kategori. Konstruksi itu dapat
	dibayangkan sebagai lokalisasi funktor dan menjadi kunci untuk memahami
	funktor turunan.

	% segment-id: o014.li2.chapter1.overview.i006
	\item Teorema funktor adjoin yang dibahas dalam
	\S~\sourcecrossref{sec:AFT}{1.11} mencirikan keberadaan adjoin kiri atau
	adjoin kanan suatu funktor di bawah syarat-syarat yang sesuai. Dalam buku
	ini, penerapan utamanya terdapat dalam
	\S~\sourcecrossref{sec:Grothendieck-cat}{2.10}, pada kajian kategori
	Grothendieck. Kegunaan teorema funktor adjoin dalam matematika tidak
	terbatas pada penerapan tersebut, sehingga bagian ini bernilai secara
	mandiri. Pengertian generator dan kogenerator
	(Definisi~\sourcecrossref{def:generators}{1.11.8}) yang dibahas di sana
	juga wajib dikuasai.
\end{itemize}

% segment-id: o014.li2.chapter1.overview.p002
Isi bab ini banyak mengambil bahan dari \cite{KS06, Rie16}.

% segment-id: o014.li2.chapter1.overview.n001
\begin{petunjukbacaan}
Teori dasar kategori abelian dan kompleks hanya memerlukan
\S\S~\sourcecrossref{sec:subquot}{1.1}--%
\sourcecrossref{sec:linear-cat}{1.4} dari bab ini.
\S\S~\sourcecrossref{sec:limit-functor}{1.5}--%
\sourcecrossref{sec:filtered-indlim}{1.6} berperan sebagai materi pendukung dan
juga akan berulang kali dirujuk dalam bab-bab lain. Teori kategori turunan
melibatkan \S\S~\sourcecrossref{sec:Kan-extension}{1.7}--%
\sourcecrossref{sec:functor-localization}{1.10}. Kegunaan materi tersebut tidak
terbatas pada topik itu, tetapi pembaca dapat menimbang kebutuhannya sendiri dan
menunda pembacaan hingga sebelum mempelajari kategori turunan. Bagian terakhir,
\S~\sourcecrossref{sec:AFT}{1.11}, hanya diterapkan dalam
\S~\sourcecrossref{sec:Grothendieck-cat}{2.10}. Namun, pengetahuan tentang
generator dan kogenerator di dalamnya bersifat wajib; pembaca dapat kembali
merujuknya bila diperlukan.
\end{petunjukbacaan}
